\section{The \texttt{Pol\_constBfield} McStas Component}
Constant magnetic field.

\subsection*{Identification}
\begin{itemize}
  \item \textbf{Site:} 
  \item \textbf{Author:} Peter Christiansen
  \item \textbf{Origin:} RISOE
  \item \textbf{Date:} July 2006
\end{itemize}

\subsection*{Description}
\begin{lstlisting}
Rectangular box with constant B field along y-axis (up). The
component can be rotated to make either a guide field or a spin
flipper. A neutron hitting outside the box opening or the box sides
is absorbed.

Example: Pol_constBfield(xwidth=0.1, yheight=0.1, zdepth=0.2, fliplambda=5.0)
\end{lstlisting}

\subsection*{Input parameters}
Parameters in \textbf{boldface} are required; the others are optional.

\begin{longtable}{p{0.22\textwidth}p{0.12\textwidth}p{0.46\textwidth}p{0.14\textwidth}}
\toprule
\textbf{Name} & \textbf{Unit} & \textbf{Description} & \textbf{Default} \\
\midrule
\endhead
\textbf{xwidth} & m & Width of opening &  \\
\textbf{yheight} & m & Height of opening &  \\
\textbf{zdepth} & m & Length of field &  \\
B & Gauss & Magnetic field along y-direction & 0 \\
fliplambda & AA & lambda for calculating B field & 0 \\
flipangle & deg & Angle flipped for lambda = fliplambda fliplambda and flipangle overrides B so a neutron with s= (1, 0, 0) and v = (0, 0, v(fliplambda)) will have s =(cos(flipangle), sin(flipangle), 0) after the section. & 180 \\
\bottomrule
\end{longtable}

\subsection*{Links}
\begin{itemize}
  \item \href{run:/home/willend/willend-McCode/mcstas-comps/optics/Pol_constBfield.comp}{Source code} for \texttt{Pol\_constBfield.comp}.
\end{itemize}
\IfFileExists{Pol_constBfield_static.tex}{\input{Pol_constBfield_static.tex}}{}